\usemodule[memos]
\usemodule[setups]

\enablemode[no-setup-main]

\setupinteraction[state=start]
\setuptextrules[style=\tf\ss,before=,after=,rulethickness=.5pt]

\setuptyping[style=\ttxx\setuplocalinterlinespace[line=2ex],
             option={context},
             before={\vbox{\hrule height .5pt\par
                     \textrule{代码示例}%
                     \hrule height .5pt\vrule width0pt depth 4pt}\par},
              after={\vbox{\hrule height .5pt\strut}},
             ]

\protected\def\colorshowtemp#1{%
  \framed[foregroundcolor=#1,width=6em,frame=off]{#1}}

\protected\def\bgcolorshowtemp#1{%
  \framed[foregroundcolor=white,background=color,backgroundcolor=#1,
           width=6em,frame=off]{#1}}

\def\showcolorrow#1{%
  \startxrow
    \startxcell {#1}                    \stopxcell
    \startxcell \colorshowtemp{soft#1}  \stopxcell
    \startxcell \colorshowtemp{light#1} \stopxcell
    \startxcell \colorshowtemp{#1}      \stopxcell
    \startxcell \colorshowtemp{dark#1}  \stopxcell
    \startxcell \colorshowtemp{deep#1}  \stopxcell
  \stopxrow}

\def\showbgcolorrow#1{%
  \startxrow
    \startxcell {#1}                      \stopxcell
    \startxcell \bgcolorshowtemp{soft#1}  \stopxcell
    \startxcell \bgcolorshowtemp{light#1} \stopxcell
    \startxcell \bgcolorshowtemp{#1}      \stopxcell
    \startxcell \bgcolorshowtemp{dark#1}  \stopxcell
    \startxcell \bgcolorshowtemp{deep#1}  \stopxcell
  \stopxrow}

\def\showpaletcolors{%
  \starttabulate[|c|c|c|c|c|]
  \NC \colorshowtemp{themecolor} \NC
      \colorshowtemp{softthemecolor} \NC
      \colorshowtemp{lightthemecolor} \NC
      \colorshowtemp{darkthemecolor} \NC
      \colorshowtemp{deepthemecolor} \NC\NR
  \NC \bgcolorshowtemp{themecolor} \NC
      \bgcolorshowtemp{softthemecolor} \NC
      \bgcolorshowtemp{lightthemecolor} \NC
      \bgcolorshowtemp{darkthemecolor} \NC
      \bgcolorshowtemp{deepthemecolor} \NC\NR
  \stoptabulate}

\starttext

\startfrontmatter
\title{memos 模块使用说明书}

\blank[big]

\placecontent

\stopfrontmatter

\startbodymatter
\chapter{模块概述}

\section{简介}

memos 是一个为 ConTeXt LMTX 设计的综合排版模块，提供了完整的书籍、文档排版解决方案。
该模块专注于中文排版，同时支持多语言混排，提供了丰富的页面布局、字体配置、颜色主题等功能。

\section{主要特性}

\startitemize[packed]
\item 多语言支持：简体中文、繁体中文、日语等
\item 丰富的字体集：Adobe、Source、macOS 等字体方案
\item 多种页面布局：moderate、kindle、ipad、kaolike 等布局样式
\item 主题颜色系统：11 种预设主题颜色调色板
\item 灵活的页眉页脚样式：book、novel、colorful 等多种风格
\item 完善的目录系统：支持多种目录样式
\item 图表浮动系统：支持文本区、边注区、全页宽等多种浮动体
\item 代码排版：支持语法高亮和行号
\item 边注系统：脚注、边注、侧注等
\item 卡片样式：Metro、Fluent、macOS、Modern、DualHeader 风格
\item 封面系统：5 种预设封面样式
\stopitemize

\section{系统要求}

该模块需要 {\bf ConTeXt LMTX (LuaMetaTeX)} 环境，不支持 MkIV 版本。
建议使用最新版本的 ConTeXt LMTX 以获得最佳体验。

\chapter{模块加载与参数}

\section{加载模块}

\starttyping
\usemodule[memos]
\stoptyping

\section{模块参数}

模块加载时可设置以下参数：

\starttabulate[|l|l|p|]
\HL
\NC 参数 \NC 默认值 \NC 说明 \NC\NR
\HL
\NC papersize \NC A4 \NC 纸张大小（A4/A5/kindle/ipad/B5/letter/legal） \NC\NR
\NC layout \NC moderate \NC 页面布局样式 \NC\NR
\NC mainlanguage \NC hans \NC 主语言 \NC\NR
\NC language \NC hans \NC 当前语言 \NC\NR
\NC lineheight \NC 1.5em \NC 行高 \NC\NR
\NC fallbackfont \NC \NC 回退字体 \NC\NR
\NC bodyfont \NC \NC 正文字体 \NC\NR
\NC fontstyle \NC rm \NC 字体风格 \NC\NR
\NC fontsize \NC 11pt \NC 字号 \NC\NR
\NC textcount \NC 40 \NC 文本区字符数 \NC\NR
\NC heightcount \NC 40 \NC 文本区行数 \NC\NR
\NC margincount \NC 0 \NC 边注区字符数 \NC\NR
\NC mode \NC \NC 运行模式 \NC\NR
\NC ratio \NC 1.414 \NC 版心比例 \NC\NR
\NC themecolor \NC cyan \NC 主题颜色 \NC\NR
\NC chapterstyle \NC default \NC 章节标题样式 \NC\NR
\NC hdrstyle \NC fctext \NC 页眉页脚样式 \NC\NR
\NC tocstyle \NC \NC 目录样式 \NC\NR
\HL
\stoptabulate

\subsection{纸张大小（papersize）}

\starttabulate[|l|p|]
\HL
\NC 选项 \NC 说明 \NC\NR
\HL
\NC A4 \NC A4 纸张（默认） \NC\NR
\NC A5 \NC A5 纸张 \NC\NR
\NC B5 \NC B5 纸张 \NC\NR
\NC letter \NC 信纸 \NC\NR
\NC legal \NC 法律用纸 \NC\NR
\NC kindle \NC Kindle 电子书尺寸 \NC\NR
\NC ipad \NC iPad 平板尺寸 \NC\NR
\HL
\stoptabulate

\subsection{页面布局（layout）}

\starttabulate[|l|p|]
\HL
\NC 选项 \NC 说明 \NC\NR
\HL
\NC moderate \NC 适中布局（默认） \NC\NR
\NC narrow \NC 窄边距布局 \NC\NR
\NC wide \NC 宽边距布局 \NC\NR
\NC extrawide \NC 超宽边距布局 \NC\NR
\NC kindle \NC Kindle 电子书布局 \NC\NR
\NC ipad \NC iPad 平板布局 \NC\NR
\NC kaolike \NC 考理科布局 \NC\NR
\NC exotic \NC 异域风格布局 \NC\NR
\NC sidemirror \NC 侧边镜像布局 \NC\NR
\NC propiorroot \NC Propiorroot 布局 \NC\NR
\NC propiorgolden \NC 黄金比例布局 \NC\NR
\HL
\stoptabulate

\subsection{语言设置（mainlanguage）}

\starttabulate[|l|p|]
\HL
\NC 选项 \NC 说明 \NC\NR
\HL
\NC hans \NC 简体中文（默认） \NC\NR
\NC hant \NC 繁体中文 \NC\NR
\NC ja \NC 日语 \NC\NR
\NC en \NC 英语 \NC\NR
\NC cn \NC 中文（简体中文别名） \NC\NR
\HL
\stoptabulate

\subsection{字体风格（fontstyle）}

\starttabulate[|l|p|]
\HL
\NC 选项 \NC 说明 \NC\NR
\HL
\NC rm \NC 罗马体（默认），适合正文 \NC\NR
\NC ss \NC 无衬线体，适合标题 \NC\NR
\NC tt \NC 等宽体，适合代码 \NC\NR
\NC hw \NC 手写体风格 \NC\NR
\HL
\stoptabulate

\subsection{运行模式（mode）}

\starttabulate[|l|p|]
\HL
\NC 选项 \NC 说明 \NC\NR
\HL
\NC print \NC 打印模式，优化输出用于打印 \NC\NR
\NC kindle \NC 电子书模式，针对电子阅读器优化 \NC\NR
\NC draft \NC 草稿模式，快速预览 \NC\NR
\NC moresize \NC 扩展字号模式，启用更多字号选项 \NC\NR
\HL
\stoptabulate

\subsection{主题颜色（themecolor）}

\starttabulate[|l|p|]
\HL
\NC 选项 \NC 说明 \NC\NR
\HL
\NC cyan \NC 青色（默认） \NC\NR
\NC red \NC 红色 \NC\NR
\NC blue \NC 蓝色 \NC\NR
\NC green \NC 绿色 \NC\NR
\NC yellow \NC 黄色 \NC\NR
\NC orange \NC 橙色 \NC\NR
\NC purple \NC 紫色 \NC\NR
\NC pink \NC 粉色 \NC\NR
\NC gray \NC 灰色 \NC\NR
\NC black \NC 黑色 \NC\NR
\NC white \NC 白色 \NC\NR
\HL
\stoptabulate

\subsection{页眉页脚样式（hdrstyle）}

\starttabulate[|l|p|]
\HL
\NC 选项 \NC 说明 \NC\NR
\HL
\NC fctext \NC 页脚居中页码（默认） \NC\NR
\NC book \NC 书籍样式 \NC\NR
\NC line \NC 书籍样式（book 别名） \NC\NR
\NC novel \NC 小说样式 \NC\NR
\NC colorful \NC 彩色样式 \NC\NR
\NC madsen \NC Madsen 样式 \NC\NR
\NC hctext \NC 页眉居中页码 \NC\NR
\NC foemargin \NC 页脚外边距页码 \NC\NR
\NC foemarginalt \NC 页脚外边距页码（靠近文本区） \NC\NR
\NC hoemargin \NC 页眉外边距样式 \NC\NR
\NC kaolike \NC 考理科风格 \NC\NR
\NC none \NC 无页眉页脚 \NC\NR
\HL
\stoptabulate

\subsection{示例}

\starttyping
% 基本用法
\usemodule[memos]

% 自定义参数
\usemodule[memos][
  papersize=A5,
  layout=moderate,
  mainlanguage=hans,
  fontsize=12pt,
  themecolor=blue,
  hdrstyle=book,
]

% 电子书模式
\usemodule[memos][mode=kindle,papersize=kindle]

% 打印模式
\usemodule[memos][mode=print]

% 草稿模式
\usemodule[memos][mode=draft]

% 更多字号选项
\usemodule[memos][mode=moresize]
\stoptyping

\chapter{多语言支持}

\section{语言设置}

memos 模块提供了完善的多语言支持，特别是对中文排版的优化。

\subsection{预设语言}

模块预设了以下语言环境：

\starttabulate[|l|l|p|]
\HL
\NC 语言代码 \NC 语言名称 \NC 说明 \NC\NR
\HL
\NC hans \NC 简体中文 \NC 使用 Adobe 简体中文字体 \NC\NR
\NC hant \NC 繁体中文 \NC 使用 Adobe 繁体中文字体 \NC\NR
\NC ja \NC 日语 \NC 使用日语字体 \NC\NR
\NC cn \NC 中文 \NC 简体中文别名 \NC\NR
\NC en \NC 英语 \NC 英语环境 \NC\NR
\HL
\stoptabulate

\subsection{切换语言}

使用 \type{\language} 命令切换语言：

\starttyping
\language[hans]  % 切换到简体中文
\language[hant]  % 切换到繁体中文
\language[ja]    % 切换到日语
\language[en]    % 切换到英语
\stoptyping

\section{自定义语言环境}

使用 \type{\defineuserlanguage} 或 \type{\setupuserlanguage} 定义新的语言环境：

\starttyping
\defineuserlanguage[mylang][
  bodyfont=adobehans,
  fallbackfont=fallback:cormorant,
  patterns={zh,en},
  script=hanzi,
  interlinespace=1.5em,
  whitespace=none,
  align={height,hanging,justified},
  indenting={first,always,2em},
]
\stoptyping

\subsection{语言参数}

\starttabulate[|l|p|]
\HL
\NC 参数 \NC 说明 \NC\NR
\HL
\NC bodyfont \NC 正文字体 \NC\NR
\NC fallbackfont \NC 回退字体 \NC\NR
\NC patterns \NC 断字模式 \NC\NR
\NC script \NC 脚本（hanzi/latin等） \NC\NR
\NC interlinespace \NC 行距 \NC\NR
\NC whitespace \NC 空白处理 \NC\NR
\NC align \NC 对齐方式 \NC\NR
\NC indenting \NC 缩进设置 \NC\NR
\HL
\stoptabulate

\chapter{字体系统}

\section{字体集}

memos 模块预定义了多种字体集，方便用户快速切换。

\subsection{Adobe 字体系列}

\starttyping
\switchtobodyfont[adobehans]   % Adobe 简体中文
\switchtobodyfont[adobehant]   % Adobe 繁体中文
\switchtobodyfont[adobenihon]  % Adobe 日语
\stoptyping

\subsection{Source 字体系列}

\starttyping
\switchtobodyfont[sourcehans]      % Source Han Sans 简体中文
\switchtobodyfont[sourcehant]      % Source Han Sans 繁体中文
\switchtobodyfont[sourcenihon]     % Source Han Sans 日语
\switchtobodyfont[sourcehanslight] % Source Han Sans Light
\stoptyping

\subsection{macOS 字体系列}

\starttyping
\switchtobodyfont[sinohans]  % PingFang SC
\switchtobodyfont[sinohant]  % PingFang TC
\switchtobodyfont[hiragino]  % Hiragino Sans
\stoptyping

\subsection{其他字体集}

\starttyping
\switchtobodyfont[genyotw]    % GenYo Gothic TW
\switchtobodyfont[genyotc]    % GenYo Gothic TC
\switchtobodyfont[shanggu]    % Shanggu
\switchtobodyfont[ipaexhans]  % IPAex 简体中文
\switchtobodyfont[ibmplexsc]  % IBM Plex SC
\stoptyping

\section{字体风格}

模块支持四种字体风格：

\starttabulate[|l|l|p|]
\HL
\NC 风格代码 \NC 风格名称 \NC 说明 \NC\NR
\HL
\NC rm \NC Roman \NC 衬线字体（默认） \NC\NR
\NC ss \NC Sans Serif \NC 无衬线字体 \NC\NR
\NC tt \NC Typewriter \NC 等宽字体 \NC\NR
\NC hw \NC Handwriting \NC 手写字体 \NC\NR
\HL
\stoptabulate

使用方法：

\starttyping
{\rm 衬线字体}
{\ss 无衬线字体}
{\tt 等宽字体}
{\hw 手写字体}
\stoptyping

\section{字号系统}

\subsection{标准字号}

\starttyping
{\tf 正常大小}
{\tfa 稍大}
{\tfb 更大}
{\tfc 很大}
{\tfd 最大}
{\tfx 稍小}
{\tfxx 更小}
\stoptyping

\subsection{扩展字号}

启用 \type{moresize} 模式后，可以使用更多字号选项：

\starttyping
\usemodule[memos][mode=moresize]

% 中文传统字号
{\o  初号}
{\i  一号}
{\ii 二号}
{\iii 三号}
{\iv 四号}
{\v  五号}

% 带小字号
{\os 小初号}
{\is 小一号}
{\iis 小二号}
{\iiis 小三号}
{\ivs 小四号}
{\vs 小五号}
\stoptyping

\chapter{页面布局}

\section{布局样式}

模块提供了多种预定义的页面布局样式：

\starttabulate[|l|p|]
\HL
\NC 布局名称 \NC 说明 \NC\NR
\HL
\NC moderate \NC 适中布局（默认），适合大多数文档 \NC\NR
\NC narrow \NC 窄边距布局，适合内容密集的文档 \NC\NR
\NC wide \NC 宽边距布局，适合边注较多的文档 \NC\NR
\NC extrawide \NC 超宽边距布局，适合大量边注 \NC\NR
\NC kindle \NC Kindle 电子书布局，适合电子阅读器 \NC\NR
\NC ipad \NC iPad 平板布局，适合平板设备 \NC\NR
\NC kaolike \NC 考理科布局，适合试卷和教材 \NC\NR
\NC exotic \NC 异域风格布局 \NC\NR
\NC sidemirror \NC 侧边镜像布局 \NC\NR
\NC propiorroot \NC Propiorroot 布局 \NC\NR
\NC propiorgolden \NC 黄金比例布局 \NC\NR
\HL
\stoptabulate

\section{布局详细参数}

每种布局都可以通过参数自定义：

\starttabulate[|l|p|]
\HL
\NC 参数 \NC 说明 \NC\NR
\HL
\NC width \NC 文本区宽度 \NC\NR
\NC height \NC 文本区高度 \NC\NR
\NC backspace \NC 左边距 \NC\NR
\NC topspace \NC 上边距 \NC\NR
\NC leftmargin \NC 左边注宽度 \NC\NR
\NC rightmargin \NC 右边注宽度 \NC\NR
\NC leftmargindistance \NC 左边注距离 \NC\NR
\NC rightmargindistance \NC 右边注距离 \NC\NR
\NC header \NC 页眉高度 \NC\NR
\NC footer \NC 页脚高度 \NC\NR
\NC headerdistance \NC 页眉距离 \NC\NR
\NC footerdistance \NC 页脚距离 \NC\NR
\NC leftedge \NC 左边缘宽度 \NC\NR
\NC rightedge \NC 右边缘宽度 \NC\NR
\NC leftedgedistance \NC 左边缘距离 \NC\NR
\NC rightedgedistance \NC 右边缘距离 \NC\NR
\NC bottom \NC 底部距离 \NC\NR
\NC bottomdistance \NC 底部距离 \NC\NR
\HL
\stoptabulate

\section{布局展示}

以下展示各布局的具体参数设置：

\doloopoverlist
  {exotic, kaolike, ipad, kindle, sidemirror,
   propiorroot, propiorgolden,
   narrow, moderate, wide, extrawide}
  {{\page\showframe\setuplayout[\recursestring]{\red\ssd \currentlayout}\showsetups\page}}

\setuplayout[default]

\section{使用方法}

\starttyping
% 使用预设布局
\usemodule[memos][layout=moderate]

% 使用 Kindle 布局
\usemodule[memos][layout=kindle,papersize=kindle]

% 使用 iPad 布局
\usemodule[memos][layout=ipad,papersize=ipad]

% 使用考理科布局
\usemodule[memos][layout=kaolike]

% 自定义布局参数
\usemodule[memos][
  textcount=42,    % 文本区宽度（字符数）
  heightcount=45,  % 文本区高度（行数）
  margincount=10,  % 边注区宽度（字符数）
  ratio=1.414,     % 版心比例（黄金比例）
]
\stoptyping

\chapter{颜色系统}

\section{基础颜色集合}

模块预定义了 11 种基础颜色，每种颜色都有 5 个层次（soft/light/normal/dark/deep）：

\subsection{文字颜色展示}

\start
\startxtable[split=yes]
\startxrow
  \startxcell 颜色 \stopxcell
  \startxcell soft \stopxcell
  \startxcell light \stopxcell
  \startxcell normal \stopxcell
  \startxcell dark \stopxcell
  \startxcell deep \stopxcell
\stopxrow
\showcolorrow{red}
\showcolorrow{blue}
\showcolorrow{yellow}
\showcolorrow{green}
\showcolorrow{cyan}
\showcolorrow{orange}
\showcolorrow{purple}
\showcolorrow{pink}
\showcolorrow{gray}
\showcolorrow{black}
\showcolorrow{white}
\stopxtable
\stop

\subsection{背景颜色展示}

\start
\startxtable[split=yes]
\startxrow
  \startxcell 颜色 \stopxcell
  \startxcell soft \stopxcell
  \startxcell light \stopxcell
  \startxcell normal \stopxcell
  \startxcell dark \stopxcell
  \startxcell deep \stopxcell
\stopxrow
\showbgcolorrow{red}
\showbgcolorrow{blue}
\showbgcolorrow{yellow}
\showbgcolorrow{green}
\showbgcolorrow{cyan}
\showbgcolorrow{orange}
\showbgcolorrow{purple}
\showbgcolorrow{pink}
\showbgcolorrow{gray}
\stopxtable
\stop

\section{主题颜色调色板}

模块预设了 11 种主题颜色调色板，每种调色板都包含完整的颜色体系：

\starttabulate[|l|l|p|]
\HL
\NC 调色板名称 \NC 主题色 \NC 说明 \NC\NR
\HL
\NC serene \NC blue \NC 宁静蓝 - 适合科技文档 \NC\NR
\NC harmony \NC green \NC 和谐绿 - 适合教育材料 \NC\NR
\NC passion \NC red \NC 热情红 - 适合重点强调 \NC\NR
\NC adventure \NC orange \NC 冒险橙 - 适合创意设计 \NC\NR
\NC optimism \NC yellow \NC 乐观黄 - 适合活泼主题 \NC\NR
\NC creativity \NC cyan \NC 创意青 - 适合现代风格（默认） \NC\NR
\NC magic \NC purple \NC 魔法紫 - 适合神秘主题 \NC\NR
\NC romance \NC pink \NC 浪漫粉 - 适合文艺作品 \NC\NR
\NC reliable \NC gray \NC 可靠灰 - 适合商务文档 \NC\NR
\NC formality \NC black \NC 正式黑 - 适合正式文档 \NC\NR
\NC innocence \NC white \NC 纯真白 - 适合简洁风格 \NC\NR
\HL
\stoptabulate

\section{调色板颜色展示}

\start
\doloopoverlist
  {serene, harmony, passion, adventure, optimism, creativity, magic, romance, reliable, formality, innocence}
  {{\red\ssd\recursestring}\par
    \setuppalet[\recursestring]\showpaletcolors}
\stop

\section{使用调色板}

\starttyping
% 方法一：通过 themecolor 参数
\usemodule[memos][themecolor=blue]

% 方法二：在文档中切换调色板
\setuppalet[serene]     % 蓝色主题
\setuppalet[harmony]    % 绿色主题
\setuppalet[passion]    % 红色主题
\setuppalet[creativity] % 青色主题
\stoptyping

\section{颜色变体}

每种主题颜色都有多个变体，形成完整的颜色体系：

\starttabulate[|l|p|]
\HL
\NC 颜色变体 \NC 说明 \NC\NR
\HL
\NC themecolor \NC 主题色 \NC\NR
\NC softthemecolor \NC 柔和主题色 \NC\NR
\NC lightthemecolor \NC 浅色主题色 \NC\NR
\NC darkthemecolor \NC 深色主题色 \NC\NR
\NC deepthemecolor \NC 极深主题色 \NC\NR
\NC softthemecolorv \NC 柔和垂直渐变色 \NC\NR
\NC optcolor \NC 选项色 \NC\NR
\NC lightoptcolor \NC 浅色选项色 \NC\NR
\HL
\stoptabulate

\section{语义颜色}

模块定义了多种语义颜色，用于不同的文档元素：

\starttabulate[|l|p|]
\HL
\NC 颜色名称 \NC 说明 \NC\NR
\HL
\NC textcolor \NC 文本颜色 \NC\NR
\NC codecolor \NC 代码颜色 \NC\NR
\NC captioncolor \NC 标题颜色 \NC\NR
\NC sidenotecolor \NC 边注颜色 \NC\NR
\NC footnotecolor \NC 脚注颜色 \NC\NR
\NC linkcolor \NC 链接颜色 \NC\NR
\NC toclinkcolor \NC 目录链接颜色 \NC\NR
\NC framecolor \NC 边框颜色 \NC\NR
\NC referencecolor \NC 引用颜色 \NC\NR
\NC titlecolor \NC 标题颜色 \NC\NR
\NC softbackground \NC 柔和背景色 \NC\NR
\NC darkbackground \NC 深色背景色 \NC\NR
\HL
\stoptabulate

\chapter{页眉页脚}

\section{页眉页脚样式}

模块提供了多种页眉页脚样式：

\starttabulate[|l|p|]
\HL
\NC 样式名称 \NC 说明 \NC\NR
\HL
\NC book \NC 书籍样式：奇数页显示章节号和标题，偶数页显示页码和章节 \NC\NR
\NC novel \NC 小说样式：页码和章节标题 \NC\NR
\NC colorful \NC 彩色样式：页码在页边彩色框中 \NC\NR
\NC madsen \NC Madsen 样式：类似书籍样式 \NC\NR
\NC hctext \NC 页眉居中页码 \NC\NR
\NC fctext \NC 页脚居中页码（默认） \NC\NR
\NC foemargin \NC 页脚外边距页码（双面打印） \NC\NR
\NC foemarginalt \NC 页脚外边距页码（靠近文本区） \NC\NR
\NC hoemargin \NC 页眉外边距样式（考理科风格） \NC\NR
\NC none \NC 无页眉页脚 \NC\NR
\HL
\stoptabulate

\section{使用方法}

\starttyping
% 书籍样式
\usemodule[memos][hdrstyle=book]

% 小说样式
\usemodule[memos][hdrstyle=novel]

% 彩色样式
\usemodule[memos][hdrstyle=colorful,themecolor=blue]

% 页脚居中页码
\usemodule[memos][hdrstyle=fctext]

% 无页眉页脚
\usemodule[memos][hdrstyle=none]
\stoptyping

\chapter{目录系统}

\section{目录样式}

\starttabulate[|l|p|]
\HL
\NC 样式名称 \NC 说明 \NC\NR
\HL
\NC default \NC 默认样式 \NC\NR
\NC kaolike \NC 考理科样式 \NC\NR
\NC none \NC 不显示目录 \NC\NR
\HL
\stoptabulate

\section{放置目录}

使用 \type{\placeTOCs} 命令放置目录：

\starttyping
% 默认目录
\placeTOCs

% 自定义目录
\placeTOCs[
  alternative=kaolike,
  title=目录,
  interlinespace=\baselineskip,
  content=yes,
  layout=moderate,
]
\stoptyping

\section{目录标题}

使用 \type{\tochead} 创建目录标题：

\starttyping
\tochead{目录}
\placecontent

\tochead{插图目录}
\placelistoffigures

\tochead{表格目录}
\placelistoftables
\stoptyping

\chapter{图表浮动}

\section{浮动体类型}

模块定义了多种浮动体类型，适应不同的排版需求：

\starttabulate[|l|l|p|]
\HL
\NC 浮动体类型 \NC 位置 \NC 说明 \NC\NR
\HL
\NC figure \NC 文本区 \NC 标准图片 \NC\NR
\NC textfig \NC 文本区 \NC 文本区图片 \NC\NR
\NC marginfig \NC 边注区 \NC 边注区图片 \NC\NR
\NC fullfig \NC 全页宽 \NC 全页宽图片 \NC\NR
\NC table \NC 文本区 \NC 标准表格 \NC\NR
\NC texttab \NC 文本区 \NC 文本区表格 \NC\NR
\NC margintab \NC 边注区 \NC 边注区表格 \NC\NR
\NC fulltab \NC 全页宽 \NC 全页宽表格 \NC\NR
\HL
\stoptabulate

\section{图片示例}

\subsection{文本区图片}

\startbuffer[textfig-demo]
\startplacetextfig[title=示例图片]
  \externalfigure[][height=4\baselineskip,width=0.6\textwidth]
\stopplacetextfig
\stopbuffer

\typebuffer[textfig-demo]

\getbuffer[textfig-demo]

\subsection{边注区图片}

\startbuffer[marginfig-demo]
\startplacemarginfig[title=边注图片]
  \externalfigure[][width=\rightmarginwidth,height=5\baselineskip]
\stopplacemarginfig
\stopbuffer

\typebuffer[marginfig-demo]

\getbuffer[marginfig-demo]

\section{表格示例}

\subsection{浮动表格}

\startbuffer[table-demo]
\startplacetable[title=示例表格]
\startxtable
  \startxrow[header]
    \startxcell 列一 \stopxcell
    \startxcell 列二 \stopxcell
    \startxcell 列三 \stopxcell
  \stopxrow
  \startxrow[body]
    \startxcell 数据 A \stopxcell
    \startxcell 数据 B \stopxcell
    \startxcell 数据 C \stopxcell
  \stopxrow
  \startxrow[lastrow]
    \startxcell 数据 D \stopxcell
    \startxcell 数据 E \stopxcell
    \startxcell 数据 F \stopxcell
  \stopxrow
\stopxtable
\stopplacetable
\stopbuffer

\typebuffer[table-demo]

\getbuffer[table-demo]

\section{图片尺寸预设}

\starttabulate[|l|p|]
\HL
\NC 预设名称 \NC 说明 \NC\NR
\HL
\NC automeasure \NC 自动尺寸 \NC\NR
\NC textmeasure \NC 文本区宽度 \NC\NR
\NC marginmeasure \NC 边注区宽度 \NC\NR
\NC fullmeasure \NC 全页宽度 \NC\NR
\HL
\stoptabulate

\chapter{代码排版}

\section{代码环境}

\subsection{基本用法}

\startbuffer[verbatim-demo]
\startverbatim
function hello() {
  print("Hello, World!")
}
\stopverbatim
\stopbuffer

\typebuffer[verbatim-demo]

\getbuffer[verbatim-demo]

\subsection{带行号的代码}

\startbuffer[lmverbatim-demo]
\startlmverbatim
function hello() {
  print("Hello, World!")
}
\stoplmverbatim
\stopbuffer

\typebuffer[lmverbatim-demo]

\getbuffer[lmverbatim-demo]

\subsection{全页宽代码}

\startbuffer[wideverbatim-demo]
\startwideverbatim
function hello() {
  print("Hello, World!")
}
\stopwideverbatim
\stopbuffer

\typebuffer[wideverbatim-demo]

\getbuffer[wideverbatim-demo]

\section{行内代码}

\startbuffer[inline-demo]
使用 \type{verbatim} 环境排版代码。
使用 \type{\textbackslash type\{...\}} 命令排版行内代码。
\stopbuffer

\typebuffer[inline-demo]

\getbuffer[inline-demo]

\chapter{注释系统}

\section{脚注}

\starttyping
这是正文内容\footnote{这是脚注内容。}
\stoptyping

\section{边注}

\starttyping
这是正文内容\sidenote{这是边注内容。}
\stoptyping

\subsection{调整边注位置}

\starttyping
\movesidenote[2\baselineskip]
这是正文内容\sidenote{这个边注会向下移动。}
\movesidenote[0em]
\stoptyping

\section{页边注}

\starttyping
这是正文内容\inoutermargin{这是页边注。}
\stoptyping

\chapter{卡片样式}

\section{基础卡片类型}

模块提供了四种基础风格的卡片：

\starttabulate[|l|l|p|]
\HL
\NC 卡片类型 \NC 风格 \NC 说明 \NC\NR
\HL
\NC MetroCard \NC Metro 风格 \NC 简洁现代，纯色背景 \NC\NR
\NC FluentCard \NC Fluent 风格 \NC 柔和渐变，强调线条 \NC\NR
\NC macOSCard \NC macOS 风格 \NC 毛玻璃效果 \NC\NR
\NC ModernCard \NC Modern 风格 \NC 渐变背景 \NC\NR
\HL
\stoptabulate

\section{基础卡片展示}

\startbuffer[metrocard-demo]
\starttextbackground[MetroCard]
  这是 Metro 风格的卡片内容。
\stoptextbackground
\stopbuffer

\startbuffer[fluentcard-demo]
\starttextbackground[FluentCard]
  这是 Fluent 风格的卡片内容。
\stoptextbackground
\stopbuffer

\startbuffer[macoscard-demo]
\starttextbackground[macOSCard]
  这是 macOS 风格的卡片内容。
\stoptextbackground
\stopbuffer

\startbuffer[moderncard-demo]
\starttextbackground[ModernCard]
  这是 Modern 风格的卡片内容。
\stoptextbackground
\stopbuffer

\subsection{Metro 风格卡片}

\getbuffer[metrocard-demo]

\subsection{Fluent 风格卡片}

\getbuffer[fluentcard-demo]

\subsection{macOS 风格卡片}

\getbuffer[macoscard-demo]

\subsection{Modern 风格卡片}

\getbuffer[moderncard-demo]

\section{DualHeader 卡片}

DualHeader 卡片是一种特殊的卡片样式，具有双栏标题框，适合展示带编号和标题的内容块。

\subsection{基本用法展示}

\startbuffer[dualheader-basic]
\startDualHeaderText[number=1,title=示例标题]
  这是 DualHeader 卡片的内容部分。
  可以包含多段文字和其他元素。
\stopDualHeaderText
\stopbuffer

\getbuffer[dualheader-basic]

\subsection{完整参数说明}

\starttabulate[|l|l|p|]
\HL
\NC 参数 \NC 默认值 \NC 说明 \NC\NR
\HL
\NC number \NC \NC 编号内容 \NC\NR
\NC title \NC \NC 标题内容 \NC\NR
\NC shownumber \NC 1 \NC 是否显示编号（0/1） \NC\NR
\NC showtitle \NC 1 \NC 是否显示标题（0/1） \NC\NR
\NC titlepos \NC 1 \NC 标题位置（0=顶部，1=左侧） \NC\NR
\NC titlealign \NC 2 \NC 标题对齐（1=左，2=中，3=右） \NC\NR
\NC numberalign \NC 2 \NC 编号对齐（1=左，2=中，3=右） \NC\NR
\NC numberbgcolor \NC lightthemecolor \NC 编号背景色 \NC\NR
\NC titlebgcolor \NC white \NC 标题背景色 \NC\NR
\NC bodybgcolor \NC white \NC 内容背景色 \NC\NR
\NC headerbgcolor \NC transparentcolor \NC 标题栏背景色 \NC\NR
\NC numbercolor \NC white \NC 编号文字颜色 \NC\NR
\NC titlecolor \NC lightthemecolor \NC 标题文字颜色 \NC\NR
\NC framecolor \NC lightthemecolor \NC 边框颜色 \NC\NR
\NC numberframecolor \NC lightthemecolor \NC 编号框边框颜色 \NC\NR
\NC titleframecolor \NC lightthemecolor \NC 标题框边框颜色 \NC\NR
\NC numberframe \NC 1 \NC 是否显示编号框边框（0/1） \NC\NR
\NC titleframe \NC 1 \NC 是否显示标题框边框（0/1） \NC\NR
\NC leftframe \NC 1 \NC 是否显示左边框（0/1） \NC\NR
\NC rightframe \NC 1 \NC 是否显示右边框（0/1） \NC\NR
\NC topframe \NC 1 \NC 是否显示上边框（0/1） \NC\NR
\NC bottomframe \NC 1 \NC 是否显示下边框（0/1） \NC\NR
\NC numberratio \NC 1 \NC 编号栏宽度比例 \NC\NR
\NC titleratio \NC 1 \NC 标题栏宽度比例 \NC\NR
\NC rulethickness \NC 1.5pt \NC 边框粗细 \NC\NR
\NC radius \NC 0pt \NC 圆角半径 \NC\NR
\NC bodybackground \NC simple \NC 内容背景类型（simple/charbox/framebox/barbox/shadowbox） \NC\NR
\NC shadowcolor \NC softgray \NC 阴影颜色（shadowbox 模式） \NC\NR
\NC symbol \NC \NC 符号 \NC\NR
\NC character \NC \NC 字符（charbox 模式） \NC\NR
\NC headerdivider \NC 0 \NC 是否显示标题分隔线（0/1） \NC\NR
\NC width \NC \type{\textwidth} \NC 卡片宽度 \NC\NR
\HL
\stoptabulate

\subsection{bodybackground 背景类型}

DualHeader 卡片支持多种内容背景类型：

\starttabulate[|l|p|]
\HL
\NC 类型 \NC 说明 \NC\NR
\HL
\NC simple \NC 简单背景，纯色填充 \NC\NR
\NC charbox \NC 字符图案背景，使用 character 参数指定的字符作为水印 \NC\NR
\NC framebox \NC 框线图案背景，点状图案 \NC\NR
\NC barbox \NC 条形背景，左侧有彩色条 \NC\NR
\NC shadowbox \NC 阴影背景，带有投影效果 \NC\NR
\HL
\stoptabulate

\subsection{进阶用法展示}

\startbuffer[dualheader-custom]
\startDualHeaderText[
  number=提示,
  title=重要说明,
  numberbgcolor=red,
  numbercolor=white,
  titlebgcolor=softred,
  titlecolor=black,
  framecolor=red,
  radius=5pt,
]
  这是一个自定义样式的 DualHeader 卡片。
  编号栏使用红色背景，标题栏使用浅红色背景。
\stopDualHeaderText
\stopbuffer

\startbuffer[dualheader-notitle]
\startDualHeaderText[shownumber=0,title=仅标题]
  这个卡片只显示标题，不显示编号。
\stopDualHeaderText
\stopbuffer

\startbuffer[dualheader-nonumber]
\startDualHeaderText[number=1,showtitle=0]
  这个卡片只显示编号，不显示标题。
\stopDualHeaderText
\stopbuffer

\startbuffer[dualheader-top]
\startDualHeaderText[titlepos=0,number=§1,title=章节说明]
  标题显示在顶部，而不是左侧。
\stopDualHeaderText
\stopbuffer

\startbuffer[dualheader-charbox]
\startDualHeaderText[
  number=★,
  title=重点内容,
  bodybackground=charbox,
  character=※,
]
  使用字符图案作为背景装饰。
\stopDualHeaderText
\stopbuffer

\startbuffer[dualheader-shadow]
\startDualHeaderText[
  number=提示,
  title=阴影效果,
  bodybackground=shadowbox,
  shadowcolor=softgray,
]
  带有阴影效果的卡片样式。
\stopDualHeaderText
\stopbuffer

\startbuffer[dualheader-barbox]
\startDualHeaderText[
  number=01,
  title=条形背景,
  bodybackground=barbox,
  leftframe=1,
]
  左侧带有彩色条的背景样式。
\stopDualHeaderText
\stopbuffer

\startbuffer[dualheader-noframe]
\startDualHeaderText[
  number=无边框,
  title=简洁样式,
  leftframe=0,
  rightframe=0,
  topframe=0,
  bottomframe=0,
]
  移除所有边框的简洁样式。
\stopDualHeaderText
\stopbuffer

\startbuffer[dualheader-radius]
\startDualHeaderText[
  number=圆角,
  title=圆角卡片,
  radius=8pt,
  numberbgcolor=blue,
  titlebgcolor=lightblue,
]
  带有圆角的卡片样式。
\stopDualHeaderText
\stopbuffer

\subsubsection{自定义颜色样式}

\getbuffer[dualheader-custom]

\subsubsection{仅标题样式}

\getbuffer[dualheader-notitle]

\subsubsection{仅编号样式}

\getbuffer[dualheader-nonumber]

\subsubsection{顶部标题样式}

\getbuffer[dualheader-top]

\subsubsection{字符图案背景}

\getbuffer[dualheader-charbox]

\subsubsection{阴影效果}

\getbuffer[dualheader-shadow]

\subsubsection{条形背景}

\getbuffer[dualheader-barbox]

\subsubsection{无边框样式}

\getbuffer[dualheader-noframe]

\subsubsection{圆角卡片}

\getbuffer[dualheader-radius]

\chapter{封面系统}

\section{封面类型}

模块提供了 5 种预设封面类型：

\starttabulate[|l|p|]
\HL
\NC 封面类型 \NC 说明 \NC\NR
\HL
\NC newbook \NC 新书封面 - 简洁现代风格 \NC\NR
\NC flyleaf \NC 扉页 - 适合书籍扉页 \NC\NR
\NC usercover \NC 用户自定义封面 - 完全可定制 \NC\NR
\NC cover \NC 标准封面 - 传统书籍封面 \NC\NR
\NC ORLY \NC O'Reilly 风格封面 - 动物图案风格 \NC\NR
\HL
\stoptabulate

\section{基本用法}

\startbuffer[cover-demo]
% 新书封面
\makecover[newbook][publisher={出版社名称}]

% 扉页
\makecover[flyleaf][publisher={出版社名称}]

% 用户自定义封面
\makecover[usercover][publisher={出版社名称}]

% 标准封面
\makecover[cover][publisher={出版社名称}]

% O'Reilly 风格封面
\makecover[ORLY][publisher={O'Reilly Media}]
\stopbuffer

\typebuffer[cover-demo]

\getbuffer[cover-demo]

\section{封面参数}

\starttabulate[|l|p|]
\HL
\NC 参数 \NC 说明 \NC\NR
\HL
\NC book \NC 书名 \NC\NR
\NC author \NC 作者 \NC\NR
\NC publisher \NC 出版社 \NC\NR
\NC series \NC 系列 \NC\NR
\NC subtitle \NC 副标题 \NC\NR
\NC version \NC 版本信息 \NC\NR
\NC animal \NC 动物图案（ORLY 封面专用） \NC\NR
\NC bookcolor \NC 书名颜色 \NC\NR
\NC authorcolor \NC 作者颜色 \NC\NR
\NC publishercolor \NC 出版社颜色 \NC\NR
\NC seriescolor \NC 系列颜色 \NC\NR
\NC bookstyle \NC 书名字体 \NC\NR
\NC authorstyle \NC 作者字体 \NC\NR
\NC publisherstyle \NC 出版社字体 \NC\NR
\NC bookalign \NC 书名对齐 \NC\NR
\NC authoralign \NC 作者对齐 \NC\NR
\NC publisheralign \NC 出版社对齐 \NC\NR
\NC seriesalign \NC 系列对齐 \NC\NR
\NC bookvspace \NC 书名垂直间距 \NC\NR
\NC authorvspace \NC 作者垂直间距 \NC\NR
\NC publishervspace \NC 出版社垂直间距 \NC\NR
\NC seriesvspace \NC 系列垂直间距 \NC\NR
\NC bookleft \NC 书名左侧装饰 \NC\NR
\NC bookright \NC 书名右侧装饰 \NC\NR
\NC afterbook \NC 书名后内容 \NC\NR
\NC lastvspace \NC 最后垂直间距 \NC\NR
\NC list \NC 内容元素列表 \NC\NR
\NC makeup \NC 页面 makeup 类型 \NC\NR
\HL
\stoptabulate

\section{各封面类型详细配置}

\subsection{newbook 封面}

新书封面是默认的封面类型，简洁现代：

\starttyping
\setupcover[newbook][
  book={书籍标题},
  author={作者姓名},
  publisher={出版社名称},
  bookcolor=themecolor,
  authorcolor=softgray,
  publishercolor=gray,
]
\makecover[newbook]
\stoptyping

\subsection{flyleaf 扉页}

扉页使用特殊的 makeup 类型，适合书籍扉页：

\starttyping
\setupcover[flyleaf][
  book={书籍标题},
  author={作者姓名},
  publisher={出版社名称},
]
\makecover[flyleaf]
\stoptyping

\subsection{usercover 用户自定义封面}

用户自定义封面提供了完整的示例配置：

\starttyping
\setupcover[usercover][
  makeup=usercover,
  list={series,subtitle,book,author,publisher},
  series={\CONTEXT\ 書籍封面示例},
  seriescolor=softgray,
  seriesvspace=,
  seriesalign=flushleft,
  subtitle={Digital Processing of Synthetic Aperture Radar Data\par
            Algorithms and Implementation},
  subtitlecolor=softgray,
  subtitlevspace=\vfill,
  subtitlealign=flushright,
  book={\hfill
        \framed[frame=on,width=8em,align=flushright,rulethickness=1pt,
                framecolor=lightoptcolor,bottomframe=off,rightframe=off,]
                {合成孔径雷达\par 成像算法与实现}},
  bookcolor=lightoptcolor,
  bookvspace=\vskip\baselineskip\relax,
  bookalign=flushright,
  author={\vbox{[加] Ian G. Cumming Frank H.Wong\quad\quad 著\par
        洪 \quad 文 \quad 胡东辉 \quad 韩 \quad 冰 \quad 等译\par
                                   吴一戎          \quad 审校\par
                                   Jyun S.Wang     \quad 排注}},
  authorcolor=softgray,
  authorvspace=\vskip2\baselineskip\relax,
  authoralign=flushright,
  publisher={这是出版社},
  publisheralign=flushright,
  publisherstyle=tt,
  publishercolor=softgray,
  publishervspace=\vskip2\baselineskip\relax,
  lastvspace=\null\vskip\baselineskip\relax,
]
\makecover[usercover]
\stoptyping

\subsection{cover 标准封面}

标准封面使用 MP 图形作为背景：

\starttyping
\setupcover[cover][
  makeup=cover,
  list={book,author,publisher,version},
  book={合成孔径雷达 \par 成像算法与实现},
  bookalign=flushleft,
  bookcolor=lightoptcolor,
  bookvspace=\vfill,
  afterbook={\blackrule[width=.5tw,height=1pt,depth=0pt,color=lightoptcolor]},
  author={Author: Soanguy},
  authorstyle=\tta,
  authorcolor=lightgray,
  authoralign=flushleft,
  authorvspace=\vskip\baselineskip\relax,
  publisher={这是出版社},
  publisherstyle=tt,
  publishercolor=lightgray,
  publishervspace=\vskip6\baselineskip\relax,
  version={version: \currentdate},
  versionstyle=tt,
  versioncolor=lightgray,
  versionalign=center,
  versionvspace=\vskip1\baselineskip\relax,
  lastvspace=,
]
\makecover[cover]
\stoptyping

\subsection{ORLY 封面}

O'Reilly 风格封面，带有动物图案：

\starttyping
\setupcover[ORLY][
  makeup=ORLY,
  list={series,animal,book,author,publisher},
  series={从入门到放弃系列权威指南},
  seriesvspace=\vskip4pt\relax,
  animal={\hfill\externalfigure[butterfly.png]},
  animalvspace=\vfil,
  bookleft={\hfill
           \startframed[background=color,backgroundcolor=themecolor,
                        frame=off,align=flushright,width=max]},
  bookright={\stopframed},
  book={合成孔径雷达 \par 成像算法与实现},
  bookalign=flushleft,
  bookcolor=white,
  bookvspace=,
  afterbook={\hfill \ttb Up and Running\par\vskip4\baselineskip\relax},
  author={{\bf O'RLY Cover} \hfill Robin Nixon},
  authorstyle=\ttb,
  authorcolor=darkgray,
  authoralign=flushleft,
  authorvspace=\vskip\baselineskip\relax,
  publisher={这是出版社},
  publisheralign=flushright,
  publisherstyle=tt,
  publishercolor=darkgray,
  publishervspace=,
  lastvspace=,
]
\makecover[ORLY]
\stoptyping

\section{封面布局}

每种封面类型都有对应的布局定义：

\starttyping
% preset 布局
\definelayout[preset]
  [backspace=15mm,rightmargin=10mm,rightedge=0mm,]

% flyleaf 布局
\definelayout[flyleaf]
  [backspace=10mm,rightmargin=10mm,rightedge=0mm,]

% usercover 布局
\definelayout[usercover][default]

% cover 布局
\definelayout[cover][default]

% ORLY 布局
\definelayout[ORLY][default][header=5mm,]
\stoptyping

\section{自定义封面}

\starttyping
% 自定义新书封面
\setupcover[newbook][
  book={我的书籍标题},
  author={作者姓名},
  publisher={出版社名称},
  bookcolor=themecolor,
  authorcolor=softgray,
  publishercolor=gray,
  bookstyle=\ssa\bf,
  authorstyle=\ss,
  publisherstyle=\tt,
  bookalign=center,
  authoralign=center,
  publisheralign=center,
]
\makecover[newbook]

% 自定义 O'Reilly 风格封面
\setupcover[ORLY][
  series={从入门到精通系列},
  book={ConTeXt 排版指南},
  author={Soanguy},
  publisher={示例出版社},
]
\makecover[ORLY]
\stoptyping

\section{封面使用建议}

\startitemize[packed]
\item 新书推荐使用 \type{newbook} 或 \type{cover} 封面
\item 技术书籍推荐使用 \type{ORLY} 封面
\item 需要完全自定义时使用 \type{usercover}
\item 扉页使用 \type{flyleaf}
\item 可以通过 \type{list} 参数控制封面元素的显示顺序
\stopitemize



\chapter{列表与枚举}

\section{有序列表}

\startbuffer[ordinals-demo]
\startordinals
  \startitem 第一项 \stopitem
  \startitem 第二项 \stopitem
  \startitem 第三项 \stopitem
\stopordinals
\stopbuffer

\typebuffer[ordinals-demo]

\getbuffer[ordinals-demo]

\section{无序列表}

\startbuffer[points-demo]
\startpoints
  \startitem 第一项 \stopitem
  \startitem 第二项 \stopitem
  \startitem 第三项 \stopitem
\stoppoints
\stopbuffer

\typebuffer[points-demo]

\getbuffer[points-demo]

\section{描述列表}

\startbuffer[description-demo]
\startdescription
  \description{术语一} 这是术语一的解释。
  \description{术语二} 这是术语二的解释。
\stopdescription
\stopbuffer

\typebuffer[description-demo]

\getbuffer[description-demo]

\section{嵌套列表}

\startbuffer[nested-demo]
\startpoints
  \startitem 第一项 \stopitem
  \startitem 
    第二项包含子列表：
    \startordinals
      \startitem 子项一 \stopitem
      \startitem 子项二 \stopitem
    \stopordinals
  \stopitem
  \startitem 第三项 \stopitem
\stoppoints
\stopbuffer

\typebuffer[nested-demo]

\getbuffer[nested-demo]

\chapter{引用与链接}

\section{引用}

\startbuffer[quote-demo]
\quote{这是引用内容。}
\stopbuffer

\startbuffer[quotation-demo]
\startquotation
  这是长引用内容。
  可以包含多段文字。
\stopquotation
\stopbuffer

\subsection{短引用}

\typebuffer[quote-demo]

\getbuffer[quote-demo]

\subsection{长引用}

\typebuffer[quotation-demo]

\getbuffer[quotation-demo]

\section{链接}

\subsection{URL 链接}

\startbuffer[url-demo]
\from[http://example.com]

\useURL[myurl][http://example.com][][示例网站]
\from[myurl]
\stopbuffer

\typebuffer[url-demo]

\getbuffer[url-demo]

\subsection{邮件链接}

\startbuffer[mail-demo]
\mailto{user@example.com}

\Mailto{user@example.com}{联系作者}
\stopbuffer

\typebuffer[mail-demo]

\getbuffer[mail-demo]

\chapter{表格样式}

\section{TABLE 表格样式}

TABLE 表格使用 \type{\startTABLE} 环境，支持以下预设样式：

\starttabulate[|l|p|]
\HL
\NC 样式名称 \NC 说明 \NC\NR
\HL
\NC trilines \NC 三线表 \NC\NR
\NC fullines \NC 全框线表格 \NC\NR
\NC alternatecolor \NC 交替行颜色 \NC\NR
\NC nonlines \NC 无框线表格 \NC\NR
\HL
\stoptabulate

\subsection{三线表（TABLE）}

\startbuffer[trilines-demo]
\startTABLE[setups=trilines]
\NC 表头一 \NC 表头二 \NC 表头三 \NC\NR
\NC 内容 A \NC 内容 B \NC 内容 C \NC\NR
\NC 内容 D \NC 内容 E \NC 内容 F \NC\NR
\stopTABLE
\stopbuffer

\getbuffer[trilines-demo]

\subsection{交替行颜色表格（TABLE）}

\startbuffer[alternate-demo]
\startTABLE[setups=alternatecolor]
\NC 表头一 \NC 表头二 \NC 表头三 \NC\NR
\NC 内容 A \NC 内容 B \NC 内容 C \NC\NR
\NC 内容 D \NC 内容 E \NC 内容 F \NC\NR
\stopTABLE
\stopbuffer

\getbuffer[alternate-demo]

\section{xtable 表格样式}

xtable 使用 \type{\startxtable} 环境，模块预设了以下行类型：

\starttabulate[|l|p|]
\HL
\NC 行类型 \NC 说明 \NC\NR
\HL
\NC header \NC 表头行，底部无框线 \NC\NR
\NC head \NC 表头行，底部无框线 \NC\NR
\NC body \NC 内容行 \NC\NR
\NC foot \NC 表尾行，顶部无框线 \NC\NR
\NC footer \NC 表尾行，顶部无框线 \NC\NR
\NC lastrow \NC 最后一行，顶部无框线 \NC\NR
\HL
\stoptabulate

\subsection{xtable 基本表格}

\startbuffer[xtable-demo]
\startxtable
  \startxrow[header]
    \startxcell 表头一 \stopxcell
    \startxcell 表头二 \stopxcell
    \startxcell 表头三 \stopxcell
  \stopxrow
  \startxrow[body]
    \startxcell 内容 A \stopxcell
    \startxcell 内容 B \stopxcell
    \startxcell 内容 C \stopxcell
  \stopxrow
  \startxrow[lastrow]
    \startxcell 内容 D \stopxcell
    \startxcell 内容 E \stopxcell
    \startxcell 内容 F \stopxcell
  \stopxrow
\stopxtable
\stopbuffer

\getbuffer[xtable-demo]

\chapter{章节样式}

\section{可用章节样式}

模块提供了多种章节标题样式，可通过 \type{chapterstyle} 参数设置：

\starttabulate[|l|p|]
\HL
\NC 样式名称 \NC 说明 \NC\NR
\HL
\NC default \NC 默认样式（默认） \NC\NR
\NC line \NC 线条样式 \NC\NR
\NC colorful \NC 彩色样式 \NC\NR
\NC madsen \NC Madsen 风格 \NC\NR
\NC rocket \NC 火箭样式 \NC\NR
\NC hexa \NC 六边形渐变样式 \NC\NR
\NC simple \NC 简洁样式 \NC\NR
\NC classics \NC 古典样式（XX 第一） \NC\NR
\NC classicnovel \NC 古典小说样式（第 X 回） \NC\NR
\NC kaolike \NC 考理科风格 \NC\NR
\NC article \NC 文章样式 \NC\NR
\NC publish \NC 出版样式 \NC\NR
\NC none \NC 无章节样式 \NC\NR
\HL
\stoptabulate

\section{使用方法}

\starttyping
\usemodule[memos][chapterstyle=madsen]
\stoptyping

\chapter{高级功能}

\section{模式控制}

\subsection{打印模式}

\starttyping
\usemodule[memos][mode=print]
\stoptyping

打印模式会优化输出用于打印，可能调整颜色和布局。

\subsection{电子书模式}

\starttyping
\usemodule[memos][mode=kindle]
\stoptyping

电子书模式针对电子阅读器优化，调整页面尺寸和布局。

\subsection{草稿模式}

\starttyping
\usemodule[memos][mode=draft]
\stoptyping

草稿模式用于快速预览，可能跳过某些耗时的排版处理。

\subsection{扩展字号模式}

\starttyping
\usemodule[memos][mode=moresize]
\stoptyping

启用更多字号选项，特别是中文传统字号。

\section{章节标题样式}

\starttyping
\usemodule[memos][chapterstyle=default]
\stoptyping

\section{组合使用}

可以组合多个模式：

\starttyping
\usemodule[memos][mode={print,moresize}]
\stoptyping

\chapter{中文数字扩展}

\section{简介}

zhnumber 模块提供了完整的中文数字转换功能，支持整数、小数、分数的中文表示，
以及日期时间、天干地支等传统中文数字应用。

\section{加载模块}

\starttyping
\usemodule[memos]
% zhnumber 功能已集成在 memos 模块中
\stoptyping

\section{整数转换}

使用 \type{\zhnumber} 命令将阿拉伯数字转换为中文数字：

\starttabulate[|l|l|]
\HL
\NC 输入 \NC 输出 \NC\NR
\HL
\NC \type{\zhnumber{0}} \NC \zhnumber{0} \NC\NR
\NC \type{\zhnumber{1}} \NC \zhnumber{1} \NC\NR
\NC \type{\zhnumber{10}} \NC \zhnumber{10} \NC\NR
\NC \type{\zhnumber{21}} \NC \zhnumber{21} \NC\NR
\NC \type{\zhnumber{100}} \NC \zhnumber{100} \NC\NR
\NC \type{\zhnumber{101}} \NC \zhnumber{101} \NC\NR
\NC \type{\zhnumber{12345}} \NC \zhnumber{12345} \NC\NR
\NC \type{\zhnumber{100000000}} \NC \zhnumber{100000000} \NC\NR
\NC \type{\zhnumber{-123}} \NC \zhnumber{-123} \NC\NR
\HL
\stoptabulate

\section{小数转换}

\starttabulate[|l|l|]
\HL
\NC 输入 \NC 输出 \NC\NR
\HL
\NC \type{\zhnumber{0.123}} \NC \zhnumber{0.123} \NC\NR
\NC \type{\zhnumber{123.45}} \NC \zhnumber{123.45} \NC\NR
\NC \type{\zhnumber{3.14159}} \NC \zhnumber{3.14159} \NC\NR
\HL
\stoptabulate

\section{分数转换}

\starttabulate[|l|l|]
\HL
\NC 输入 \NC 输出 \NC\NR
\HL
\NC \type{\zhnumber{1/2}} \NC \zhnumber{1/2} \NC\NR
\NC \type{\zhnumber{1/3}} \NC \zhnumber{1/3} \NC\NR
\NC \type{\zhnumber{3/4}} \NC \zhnumber{3/4} \NC\NR
\HL
\stoptabulate

\section{日期时间}

\subsection{日期转换}

\starttabulate[|l|l|]
\HL
\NC 命令 \NC 输出 \NC\NR
\HL
\NC \type{\zhdate{2024/1/1}} \NC \zhdate{2024/1/1} \NC\NR
\NC \type{\zhdate{2024/12/31}} \NC \zhdate{2024/12/31} \NC\NR
\NC \type{\zhtoday} \NC \zhtoday \NC\NR
\HL
\stoptabulate

\subsection{时间转换}

\starttabulate[|l|l|]
\HL
\NC 命令 \NC 输出 \NC\NR
\HL
\NC \type{\zhtime{0:00}} \NC \zhtime{0:00} \NC\NR
\NC \type{\zhtime{14:30}} \NC \zhtime{14:30} \NC\NR
\NC \type{\zhcurrtime} \NC \zhcurrtime \NC\NR
\HL
\stoptabulate

\subsection{星期转换}

\starttabulate[|l|l|]
\HL
\NC 命令 \NC 输出 \NC\NR
\HL
\NC \type{\zhweekday{2024/1/1}} \NC \zhweekday{2024/1/1} \NC\NR
\HL
\stoptabulate

\section{天干地支}

\subsection{天干}

\starttabulate[|l|l|]
\HL
\NC 命令 \NC 输出 \NC\NR
\HL
\NC \type{\zhtiangan{1}} \NC \zhtiangan{1} \NC\NR
\NC \type{\zhtiangan{2}} \NC \zhtiangan{2} \NC\NR
\NC \type{\zhtiangan{10}} \NC \zhtiangan{10} \NC\NR
\HL
\stoptabulate

完整天干：\dorecurse{10}{\zhtiangan{#1}\ifnum#1<10，\fi}

\subsection{地支}

\starttabulate[|l|l|]
\HL
\NC 命令 \NC 输出 \NC\NR
\HL
\NC \type{\zhdizhi{1}} \NC \zhdizhi{1} \NC\NR
\NC \type{\zhdizhi{2}} \NC \zhdizhi{2} \NC\NR
\NC \type{\zhdizhi{12}} \NC \zhdizhi{12} \NC\NR
\HL
\stoptabulate

完整地支：\dorecurse{12}{\zhdizhi{#1}\ifnum#1<12，\fi}

\subsection{干支}

\starttabulate[|l|l|]
\HL
\NC 命令 \NC 输出 \NC\NR
\HL
\NC \type{\zhganzhi{1}} \NC \zhganzhi{1} \NC\NR
\NC \type{\zhganzhi{60}} \NC \zhganzhi{60} \NC\NR
\NC \type{\zhganzhinian{2024}} \NC \zhganzhinian{2024} \NC\NR
\HL
\stoptabulate

\section{数字样式}

支持三种数字样式：

\starttabulate[|l|l|l|]
\HL
\NC 样式 \NC 说明 \NC 示例 \NC\NR
\HL
\NC normal \NC 普通中文数字（默认） \NC \zhnumber[normal]{12345} \NC\NR
\NC cap \NC 大写中文数字（财务用） \NC \zhnumber[cap]{12345} \NC\NR
\NC all \NC 扩展中文数字（含廿、卅、卌） \NC \zhnumber[all]{23} \NC\NR
\HL
\stoptabulate

\starttyping
\zhnumber[normal]{12345}  % 一万二千三百四十五
\zhnumber[cap]{12345}     % 壹万贰仟叁佰肆拾伍
\zhnumber[all]{23}        % 廿三
\stoptyping

\section{零的设置}

默认使用「〇」表示零，可以切换为「零」：

\starttyping
% 使用〇（默认）
\setupzhnumber[zero=〇]
\zhnumber{101}  % 一百〇一

% 使用零
\setupzhnumber[zero=零]
\zhnumber{101}  % 一百零一
\stoptyping

\section{转换器}

zhnumber 注册了以下转换器，可在 \type{\convertnumber} 中使用：

\starttabulate[|l|l|]
\HL
\NC 转换器 \NC 说明 \NC\NR
\HL
\NC tiangan \NC 天干（1-10） \NC\NR
\NC dizhi \NC 地支（1-12） \NC\NR
\NC ganzhi \NC 干支（1-60） \NC\NR
\HL
\stoptabulate

\starttyping
\convertnumber{tiangan}{1}  % 甲
\convertnumber{dizhi}{1}    % 子
\convertnumber{ganzhi}{1}   % 甲子
\stoptyping

\chapter{中文索引排序}

\section{简介}

zhindex 模块提供了中文索引排序功能，支持拼音排序、字母排序和笔画排序三种方式。

\section{加载模块}

\starttyping
\usemodule[memos]
\usemodule[zhindex]
\stoptyping

\section{排序方法}

\starttabulate[|l|p|]
\HL
\NC 排序方法 \NC 说明 \NC\NR
\HL
\NC zh-pinyin \NC 按拼音排序，同音字按拼音字母顺序排列 \NC\NR
\NC zh-alpha \NC 按字母顺序排序，汉字按拼音首字母分组 \NC\NR
\NC zh-stroke \NC 按笔画数排序，同笔画数按笔顺代码排列 \NC\NR
\HL
\stoptabulate

\section{使用方法}

\subsection{拼音排序}

\starttyping
\setupregister[index][
  n=3,
  alternative=A,
  language=zh-pinyin,
]

\placeindex
\stoptyping

\subsection{字母排序}

\starttyping
\setupregister[index][
  n=3,
  alternative=A,
  language=zh-alpha,
]

\placeindex
\stoptyping

\subsection{笔画排序}

\starttyping
\setupregister[index][
  n=3,
  alternative=A,
  language=zh-stroke,
]

\placeindex
\stoptyping

\section{添加索引条目}

使用 \type{\index} 命令添加索引条目：

\starttyping
\index{北京}
\index{上海}
\index{广州}

% 带子条目
\index{水果+苹果}
\index{水果+香蕉}
\stoptyping

\section{数据文件}

zhindex 模块需要以下数据文件：

\starttabulate[|l|p|]
\HL
\NC 文件 \NC 说明 \NC\NR
\HL
\NC pinyin.txt \NC 汉字拼音数据 \NC\NR
\NC sunwb_strokeorder.txt \NC 汉字笔画数据 \NC\NR
\HL
\stoptabulate

\section{排序特点}

\startitemize[packed]
\item 支持中西文混合排序
\item 数字排在最前，字母排在最后
\item 同类内按各自规则排序
\item 笔画排序支持按笔顺细分
\stopitemize

\chapter{最佳实践}

\section{文档结构}

推荐的文档结构：

\starttyping
\usemodule[memos][
  papersize=A4,
  layout=moderate,
  mainlanguage=hans,
  fontsize=11pt,
  themecolor=cyan,
  hdrstyle=fctext,
]

\starttext

\startfrontmatter
  \makecover[cover][publisher={出版社}]
  \placeTOCs
  \title{前言}
  前言内容...
\stopfrontmatter

\startbodymatter
  \part{第一部分}
  \chapter{第一章}
  章节内容...
\stopbodymatter

\startbackmatter
  \chapter{附录}
  附录内容...
\stopbackmatter

\stoptext
\stoptyping

\section{字体选择建议}

\startitemize[packed]
\item 正式文档：使用 Adobe 字体系列（adobehans/adobehant）
\item 开源项目：使用 Source 字体系列（sourcehans/sourcehant）
\item macOS 用户：可以使用系统字体（sinohans/sinohant）
\item 需要轻量字体：使用 Source Han Sans Light（sourcehanslight）
\stopitemize

\section{颜色搭配建议}

\startitemize[packed]
\item 科技文档：serene（蓝色）、creativity（青色）
\item 文学作品：formality（黑色）、innocence（白色）
\item 教育材料：harmony（绿色）、optimism（黄色）
\item 商务文档：reliable（灰色）、formality（黑色）
\item 创意设计：adventure（橙色）、magic（紫色）、romance（粉色）
\item 重点强调：passion（红色）
\stopitemize

\section{布局选择建议}

\startitemize[packed]
\item 一般文档：moderate（适中布局）
\item 电子书：kindle（Kindle 布局）
\item 平板阅读：ipad（iPad 布局）
\item 试卷教材：kaolike（考理科布局）
\stopitemize

\chapter{故障排除}

\section{常见问题}

\subsection{字体未找到}

确保已安装相应的字体，或切换到已安装的字体集。

\subsection{页面布局异常}

检查 \type{textcount}、\type{heightcount}、\type{margincount} 参数是否合理。

\subsection{颜色显示不正确}

确保 \type{themecolor} 参数值在预设颜色列表中，或使用 \type{\setuppalet} 切换调色板。

\subsection{目录不显示}

检查是否使用了 \type{\placeTOCs} 或 \type{\placecontent} 命令。

\subsection{封面显示异常}

检查封面参数是否正确设置，特别是 \type{book}、\type{author}、\type{publisher} 等必需参数。

\section{调试技巧}

\starttyping
% 启用调试模式
\usemodule[memos][mode=draft]

% 显示页面布局
\showframe

% 显示网格
\showgrid
\stoptyping
\stopbodymatter

\startappendices

\chapter{完整示例}

\startbuffer[complete-example]
\usemodule[memos][
  papersize=A4,
  layout=moderate,
  mainlanguage=hans,
  fontsize=11pt,
  themecolor=blue,
  hdrstyle=book,
  mode=moresize,
]

\starttext

\startfrontmatter
  \makecover[cover][publisher={示例出版社}]
  \placeTOCs
\stopfrontmatter

\startbodymatter
  \part{第一部分：基础篇}
  
  \chapter{简介}
  
  \section{背景}
  这是背景介绍\footnote{这是一个脚注}。
  
  \section{目的}
  这是目的说明\sidenote{这是一个边注}。
  
  \startplacefigure[title=示例图片]
    \externalfigure[][width=\textwidth,height=3\baselineskip]
  \stopplacefigure
  
  \startDualHeaderText[number=1,title=重要提示]
    这是一个 DualHeader 卡片示例。
    可以用于突出显示重要内容。
  \stopDualHeaderText
  
  \starttextbackground[MetroCard]
    这是一个 Metro 风格的卡片。
  \stoptextbackground
  
  \startverbatim
  function example() {
    return "Hello, World!";
  }
  \stopverbatim
  
\stopbodymatter

\startbackmatter
  \chapter{附录}
  附录内容。
\stopbackmatter

\stoptext
\stopbuffer

\typebuffer[complete-example]

\stopappendices


\stoptext
